\documentclass[11pt,a4paper]{article}
\usepackage[margin=1in]{geometry}
\usepackage[ngerman]{babel}
\usepackage[utf8]{inputenc}
\usepackage{graphicx}
\usepackage{amsmath}
\usepackage{array}
\usepackage{booktabs}
\usepackage{multicol}
\usepackage{enumitem}
\usepackage{tikz}

\newcommand{\emptybox}{%
\begin{tikzpicture}[scale=1, baseline=0.2cm]
\draw (0,0)rectangle(3,1);
\end{tikzpicture}
}


\parindent=0pt

\title{Aktivität 2: Schlüssel mit realistischeren grossen Zahlen erzeugen}
\author{}
\date{}

\begin{document}

\maketitle

Im echten Leben verwendet die Verschlüsselung viel grössere Zahlen, die sehr schwer zu erraten wären. In Aktivität 2 werden wir realistischere grössere Zahlen verwenden.

Das Erstellen von Modulartabellen würde sehr lange dauern, daher werden wir eine spezielle mathematische Technik verwenden, eine Art Abkürzung zum Finden modularer Inverser.

Die zweite „Nachricht" wird eine ganze Zahl sein, die einem Drei-Buchstaben-Wort entspricht, das Sie aus der beigefügten Tabelle auswählen werden.

\section*{Wählen Sie Ihre Schlüssel}

\begin{enumerate}[itemsep=.5cm]
    \item Wählen Sie eine Zahl $A$ zwischen 25 und 50. \hfill $A = $ \emptybox

    \item Wählen Sie eine zweite Zahl $B$ zwischen 25 und 50. \hfill $B = $ \emptybox

    \item Multiplizieren Sie $A$ und $B$. Nennen Sie diese Zahl $C$. \hfill $C = $ \emptybox

    \item Subtrahieren Sie 1. Nennen Sie diese Zahl $D$. \hfill $D = $ \emptybox
\end{enumerate}

\begin{itemize}[itemsep=0.5cm]
    \item Ihr \textbf{öffentlicher Schlüssel} ist $A + B + C$ \hfill Öffentlicher Schlüssel $= $ \emptybox

    \item Ihr \textbf{Verschlüsselungsschlüssel} ist $A + D$. \\ \phantom{text}\hfill Verschlüsselungsschlüssel $= $ \emptybox

    \item Ihr \textbf{Entschlüsselungs- oder privater Schlüssel} ist $B + D$. 
    \\\phantom{.}\hfill Entschlüsselungsschlüssel $= $ \emptybox
\end{itemize}

\textbf{Teilen Sie allen anderen ihren öffentlichen und Verschlüsselungsschlüssel mit. Ihr Entschlüsselungsschlüssel bleibt geheim!}

Wenn Sie sich nun die Mühe machen würden, es zu überprüfen, würden Sie feststellen, dass:

\begin{center}
\textbf{Verschlüsselungsschlüssel $\times$ Entschlüsselungsschlüssel $= 1$ im Modulus unseres öffentlichen Schlüssels.}
\end{center}

Wir sind endlich bereit, unsere geheime Nachricht zu \textbf{verschlüsseln} und zu \textbf{entschlüsseln}, da wir zwei Zahlen gefunden haben, deren Produkt 1 im Modulus unseres privaten Schlüssels ist.

\newpage

\section*{Verschlüsselung}

Wählen Sie eine Nachricht zur Übertragung aus dem beigefügten Blatt mit Drei-Buchstaben-Wörtern.

\begin{enumerate}
    \setcounter{enumi}{4}
    \item Notieren Sie die Zahl, die Ihrem Wort in der Tabelle entspricht, und multiplizieren Sie sie mit Ihrem \textbf{Verschlüsselungsschlüssel}

    \[
    F = \text{Nachricht (aus Tabelle)} \times \text{Verschlüsselungsschlüssel} = \emptybox
    \]

    \item Teilen Sie $F$ durch den \textbf{öffentlichen Schlüssel} und schreiben Sie nur die Ziffern vor dem Dezimalpunkt auf. \hfill $G = $ \emptybox

    \item Multiplizieren Sie $G$ mit Ihrem \textbf{öffentlichen Schlüssel}. \hfill $H = $ \emptybox

    \item Subtrahieren Sie $H$ vom Ergebnis in $F$ oben. \hfill $I = $ \emptybox

    Das Ergebnis ist Ihr verschlüsseltes geheimes Wort. \hfill \textbf{Verschlüsselte Nachricht}
\end{enumerate}

\section*{Entschlüsselung}

\begin{enumerate}
    \setcounter{enumi}{8}
    \item Multiplizieren Sie die verschlüsselte Nachricht in $I$ oben mit dem \textbf{Entschlüsselungsschlüssel}. \\ \phantom{text}\hfill $J = $ \emptybox

    \item Teilen Sie das Ergebnis in $J$ durch den \textbf{öffentlichen Schlüssel}. Schreiben Sie nur die Ziffern vor dem Dezimalpunkt auf. \hfill $K = $ \emptybox

    \item Multiplizieren Sie das Ergebnis in $K$ mit dem \textbf{öffentlichen Schlüssel}.\\ \phantom{text} \hfill $L = $ \emptybox

    \item Subtrahieren Sie das Ergebnis in $L$ vom Ergebnis in $J$. \hfill $M = $ \emptybox
\end{enumerate}

Das Ergebnis ist Ihre entschlüsselte Nachricht -- eine Zahl, die Ihr geheimes Drei-Buchstaben-Wort enthüllt, das nur denjenigen bekannt ist, die Ihren privaten \textbf{Entschlüsselungsschlüssel} zusammen mit Ihrem \textbf{öffentlichen Schlüssel} haben.

\newpage

\section*{Tabelle der Drei-Buchstaben-Wörter}

\begin{multicols}{6}
\small
\begin{tabular}{cl}
\toprule
\# & Wort \\
\midrule
1 & ace \\
2 & act \\
3 & add \\
4 & ado \\
5 & ads \\
6 & adz \\
7 & aft \\
8 & age \\
9 & ago \\
10 & aha \\
11 & aid \\
12 & ail \\
13 & aim \\
14 & air \\
15 & alb \\
16 & ale \\
17 & all \\
18 & amp \\
19 & and \\
20 & ani \\
21 & ant \\
22 & any \\
23 & ape \\
24 & apt \\
25 & arc \\
26 & are \\
27 & ark \\
28 & arm \\
29 & art \\
30 & ash \\
31 & ask \\
32 & asp \\
33 & ate \\
34 & auk \\
35 & awe \\
36 & awl \\
37 & axe \\
38 & aye \\
39 & baa \\
40 & bad \\
41 & bag \\
42 & bah \\
43 & ban \\
44 & bar \\
45 & bat \\
46 & bay \\
47 & bed \\
48 & bee \\
\bottomrule
\end{tabular}

\columnbreak

\begin{tabular}{cl}
\toprule
\# & Wort \\
\midrule
49 & beg \\
50 & bet \\
51 & bib \\
52 & bid \\
53 & big \\
54 & bin \\
55 & bit \\
56 & boa \\
57 & bob \\
58 & bog \\
59 & boo \\
60 & bop \\
61 & bow \\
62 & box \\
63 & boy \\
64 & brr \\
65 & bud \\
66 & bug \\
67 & bum \\
68 & bun \\
69 & bur \\
70 & bus \\
71 & but \\
72 & buy \\
73 & bye \\
74 & cab \\
75 & cad \\
76 & cam \\
77 & can \\
78 & cap \\
79 & car \\
80 & cat \\
81 & caw \\
82 & chi \\
83 & cob \\
84 & cod \\
85 & cog \\
86 & con \\
87 & coo \\
88 & cop \\
89 & cot \\
90 & cow \\
91 & cox \\
92 & coy \\
93 & cry \\
94 & cub \\
95 & cud \\
96 & cue \\
\bottomrule
\end{tabular}

\columnbreak

\begin{tabular}{cl}
\toprule
\# & Wort \\
\midrule
97 & cup \\
98 & cur \\
99 & cut \\
100 & dab \\
101 & dad \\
102 & dam \\
103 & day \\
104 & deb \\
105 & den \\
106 & dew \\
107 & did \\
108 & die \\
109 & dig \\
110 & dim \\
111 & din \\
112 & dip \\
113 & dis \\
114 & doc \\
115 & doe \\
116 & dog \\
117 & don \\
118 & dos \\
119 & dot \\
120 & dry \\
121 & dub \\
122 & dud \\
123 & due \\
124 & dug \\
125 & duh \\
126 & dun \\
127 & duo \\
128 & dye \\
129 & ear \\
130 & eat \\
131 & ebb \\
132 & eel \\
133 & egg \\
134 & ego \\
135 & eke \\
136 & elf \\
137 & elk \\
138 & ell \\
139 & elm \\
140 & ems \\
141 & emu \\
142 & end \\
143 & eon \\
144 & era \\
\bottomrule
\end{tabular}

\columnbreak

\begin{tabular}{cl}
\toprule
\# & Wort \\
\midrule
145 & ere \\
146 & erg \\
147 & err \\
148 & eta \\
149 & eve \\
150 & ewe \\
151 & eye \\
152 & fad \\
153 & fan \\
154 & far \\
155 & fat \\
156 & fax \\
157 & fed \\
158 & fee \\
159 & fen \\
160 & fer \\
161 & few \\
162 & fey \\
163 & fez \\
164 & fib \\
165 & fie \\
166 & fig \\
167 & fin \\
168 & fir \\
169 & fit \\
170 & fix \\
171 & flu \\
172 & fly \\
173 & fob \\
174 & foe \\
175 & fog \\
176 & fop \\
177 & for \\
178 & fox \\
179 & fro \\
180 & fry \\
181 & fun \\
182 & fur \\
183 & gab \\
184 & gad \\
185 & gag \\
186 & gal \\
187 & gap \\
188 & gas \\
189 & gee \\
190 & gel \\
191 & gem \\
192 & get \\
\bottomrule
\end{tabular}

\columnbreak

\begin{tabular}{cl}
\toprule
\# & Wort \\
\midrule
193 & gig \\
194 & gin \\
195 & gnu \\
196 & gob \\
197 & god \\
198 & goo \\
199 & gos \\
200 & got \\
201 & gum \\
202 & gun \\
203 & gut \\
204 & guy \\
205 & gym \\
206 & gyp \\
207 & had \\
208 & hag \\
209 & hah \\
210 & ham \\
211 & has \\
212 & hat \\
213 & haw \\
214 & hay \\
215 & hem \\
216 & hen \\
217 & hep \\
218 & her \\
219 & hes \\
220 & hew \\
221 & hex \\
222 & hey \\
223 & hid \\
224 & hie \\
225 & him \\
226 & hip \\
227 & his \\
228 & hit \\
229 & hob \\
230 & hod \\
231 & hoe \\
232 & hog \\
233 & hop \\
234 & hos \\
235 & hot \\
236 & how \\
237 & hub \\
238 & hue \\
239 & hug \\
240 & huh \\
\bottomrule
\end{tabular}

\columnbreak

\begin{tabular}{cl}
\toprule
\# & Wort \\
\midrule
241 & hum \\
242 & hut \\
243 & ice \\
244 & icy \\
245 & ids \\
246 & ifs \\
247 & ilk \\
248 & ill \\
249 & imp \\
250 & ink \\
251 & inn \\
252 & ins \\
253 & ion \\
254 & ire \\
255 & irk \\
256 & ism \\
257 & its \\
258 & ivy \\
259 & jab \\
260 & jag \\
261 & jam \\
262 & jar \\
263 & jaw \\
264 & jay \\
265 & jet \\
266 & jib \\
267 & jig \\
268 & job \\
269 & jog \\
270 & jot \\
271 & joy \\
272 & jug \\
273 & jut \\
274 & keg \\
275 & ken \\
276 & key \\
277 & kid \\
278 & kin \\
279 & kit \\
280 & lab \\
281 & lad \\
282 & lag \\
283 & lam \\
284 & lap \\
285 & law \\
286 & lax \\
287 & lay \\
288 & lea \\
\bottomrule
\end{tabular}
\end{multicols}

\vfill

\noindent\small
\textcopyright{} 2016 Education Services Australia Ltd, sofern nicht anders angegeben.\\
Creative Commons BY 4.0 Lizenz, sofern nicht anders angegeben.

\end{document}
